\documentclass[11pt]{article}
\usepackage{amsmath,amssymb,amsthm}
\newtheorem{theorem}{Theorem}
\newtheorem{proposition}{Proposition}
\title{Local reciprocal zero branches from sum histograms}
\author{}\date{}
\begin{document}\maketitle
\paragraph{Status.}
The branch identity, counting, regularity, and cube admissibility below
are complete written arguments. The companion program checks small
instances with exact rational arithmetic. There is no asymptotic sign
coverage theorem here, and these results do not establish either requested
bound for worst-case head complexity.

\begin{theorem}
Fix $H\geq1$ and $k\geq1$, and permutations $\rho_1,\ldots,\rho_k$
of $[H]$, with $\rho_1(h)=h$. Set $\rho_0(h)=0$. There is a family
$\mathcal F\subseteq S_H$ of size
\[
 |\mathcal F|\geq H!\,2^{-3kH}
\]
such that, for every positive $\epsilon$ and every $l=0,\ldots,k$,
the rational function
\[
 S_{\pi,l}(t)=\sum_{h=1}^H
       \frac1{t-\epsilon(\pi(h)+\rho_l(h))}
\]
is independent of $\pi\in\mathcal F$ wherever defined.
\end{theorem}
\begin{proof}
For $l\geq1$, record the histogram of $\pi(h)+\rho_l(h)$ on
$\{2,\ldots,2H\}$. Its entries are nonnegative integers summing to
$H$, so there are at most $\binom{3H-2}{H}\leq2^{3H}$ possibilities.
Pigeonhole the $k$ histograms simultaneously. For $l=0$ the histogram
is always that of $[H]$. Equal histograms give the displayed rational
identity by grouping equal denominators.
\end{proof}

For example, if $k=\lfloor(\log_2H)/12\rfloor\geq1$, then
$H!\geq(H/2)^{H/2}$ gives
\[
 \log_2|\mathcal F|
 \geq \tfrac14H\log_2H-\tfrac12H
 \geq\tfrac18H\log_2H \qquad(H\geq16).
\]
Thus $k$ of order $\log H$ is compatible with exponentially many
branches on the $H\log H$ scale. This count concerns representations,
not signs.

\begin{proposition}
Choose $2H$ distinct exterior queries $t$ in group $0$, and at least
one exterior query in group $1$. At every branch $\pi$, the score map
\[
 (a,p)\longmapsto
 \left(\sum_h\frac{a_h}{t-p_h-\epsilon\rho_l(h)}\right)_{(l,t)}
\]
has derivative of rank $2H$ at $a_h=1$, $p_h=\epsilon\pi(h)$.
Moreover the full derivative images at two different branches are
distinct subspaces of the query space.
\end{proposition}
\begin{proof}
The derivative is
\[
 J_\pi(\dot a,\dot p)_{l,t}
 =\sum_h\left(
 \frac{\dot a_h}{t-\epsilon(\pi(h)+\rho_l(h))}
 +\frac{\dot p_h}{(t-\epsilon(\pi(h)+\rho_l(h)))^2}\right).
\]
On group $0$, a vanishing vector clears to a polynomial of degree at
most $2H-1$ with $2H$ distinct zeros. Hence that polynomial vanishes
identically; uniqueness of partial fractions gives $\dot a=\dot p=0$.

Suppose two derivative images coincide. The group-$0$ restriction
uniquely identifies coefficients at each pole. Thus the direction
$\dot a_j=1$ at branch $\pi$ must match the direction
$\dot a_i=1$ at branch $\tau$, where $\tau(i)=\pi(j)$, with every
other numerator and pole direction zero. At the group-$1$ query,
equality gives
\[
 \frac1{t-\epsilon(\pi(j)+j)}
 =\frac1{t-\epsilon(\pi(j)+i)}.
\]
Therefore $i=j$ for every $j$, and $\tau=\pi$.
\end{proof}

The group-$1$ head labels also show that these branches are genuinely
different multisets of denominator profiles. They are not just a
reordering of the same heads. The common score is a regular value of
each local $2H$-parameter sheet after projection onto group $0$.
This does not assert that it is a regular value of any larger family.

\paragraph{Signed column numerators.}
One may fix signs $b_{l,h}\in\{-1,1\}$ in each group and require
the positive and negative histograms separately equal. There are at
most $\binom{5H-1}{H}\leq2^{5H}$ signed histograms, since $H$ tokens
are allocated among at most $4H$ signed bins. Consequently a common
score family for
\[
 \sum_h\frac{b_{l,h}}{t-\epsilon(\pi(h)+\rho_l(h))}
\]
has size at least $H!2^{-5kH}$. Keep $b_{0,h}=b_{1,h}=1$ to preserve
the preceding rank and distinctness proof. The derivative pole
directions in other groups acquire the factor $b_{l,h}$, while the
row-numerator directions do not. This extension permits independently
chosen signed masks in the groups; it still gives no orthant count.

\paragraph{Strict full-cube denominators.}
Let there be $r$ local rows and $q$ queries, each query $j$ belonging
to a group $l(j)$. Write $t_j=1-\eta_j$, where
$0<\eta_j<1/(8q)$. Choose one branch $\pi_i$ per row and
\[
 0<\epsilon<\frac1{16H(r+q)}.
\]
Then
\[
 B_h(z,w)=1-\epsilon\sum_i\pi_i(h)z_i
       -\sum_j(\eta_j+\epsilon\rho_{l(j)}(h))w_j
\]
has strictly negative slopes and is positive on the entire Boolean
$(r+q)$-cube. Indeed its total slope magnitude is less than
$1/16+1/8+1/16=1/4$, so its value is greater than $3/4$.
Its grid value at $(e_i,e_j)$ is exactly
$t_j-\epsilon(\pi_i(h)+\rho_{l(j)}(h))$.
Small independent perturbations of the row pole slopes retain these
strict inequalities. The affine numerator
$\sum_i\dot a_{hi}z_i+\sum_jb_{l(j),h}w_j$ supplies the independent
row numerator perturbations at the same grid.

A fixed reference branch can be subtracted with another $H$ heads,
using the same admissible construction with that branch at every row.
Thus every exact common-score cancellation here can be implemented
with $2H$ strict grid heads. The existing checked general axis repair
then has linear cost. No saving is claimed from this observation.

\paragraph{Small exact sign diagnostic.}
For $H=3$, $\rho_1=\mathrm{id}$, the fiber
$(2,3,1),(3,1,2)$ has identical histograms. At the compact grids in the
checker, its two six-dimensional tangent spaces have exactly the same
$126$ of $128$ orthants on seven queries. Distinct tangent spaces
therefore need not contribute different signs.

For $H=4$, the fiber
\[
 (3,4,2,1),\quad(4,2,3,1),\quad(4,3,1,2)
\]
has common sum histogram $(4,5,5,6)$. The checker gives nine exterior
queries for which two rank-eight derivative images have normal vectors
with distinct sign words and no zero entries. A codimension-one
subspace with such normal $v$ misses precisely the two open orthants
$\pm\operatorname{sign}(v)$: every other orthant contains both signs
among the summands of $v\cdot x$, and their magnitudes can be balanced.
The union of these two derivative images therefore meets all $512$
open orthants. This exact finite result is checked at three positive
values of $\epsilon$; no sufficiently-small-$\epsilon$ statement is
inferred from these three samples.

\paragraph{The unresolved step.}
Taking $2H$ queries in each group gives $q=2H(k+1)$, of order
$H\log H$. Neither the large fiber nor distinct tangent images
establishes that their union has $\exp(cH\log H)$ sign patterns, let
alone all orthants or a uniform positive fraction. These are the
missing statements. Taking all $\rho_l$ equal already shows why the
branch-count argument alone cannot work: all query functions then
belong to a fixed space of dimension $O(H)$, irrespective of $k$.
Recovering an individual histogram bin by a linear combination of
many exterior queries is not a free operation in the local compiler;
using $O(H)$ resolvents for each bin would lose the proposed saving.
\end{document}
